\chapter{Disability Services and the NDIS}

\section*{Definition and Overview}
The National Disability Insurance Scheme (NDIS) is Australia's national scheme that funds supports for people with permanent and significant disability. The NDIS provides funding for reasonable and necessary supports that help participants with daily living, therapy, equipment, and community participation. It is not a welfare payment but an insurance scheme designed to give people with disability choice and control over their supports. This chapter explains how the NDIS works, who is eligible, and how to access disability services through it.

\section*{Epidemiology}
The NDIS supports over 500,000 Australians with disability, including people with physical, intellectual, sensory, psychosocial, and developmental disabilities. Around one in six Australians has some form of disability. The scheme has grown substantially since its introduction, and it funds a wide range of supports from allied health therapy to daily personal care. Many people with disability also use mainstream health, education, and community services alongside their NDIS supports.

\section*{Aetiology and Pathophysiology}
\begin{itemize}
    \item \textbf{NDIS purpose:} to fund reasonable and necessary supports related to a person's disability
    \item \textbf{Eligibility:} permanent and significant disability, assessed against access criteria including age (under 65) and residency
    \item \textbf{Participant plans:} individualised plans funded through budgets for core supports, capacity building, and capital supports
    \item \textbf{Core supports:} daily living, personal care, community participation, and consumables
    \item \textbf{Capacity building:} therapies, training, and support to build skills and independence
    \item \textbf{Capital supports:} assistive technology and home modifications
    \item \textbf{Providers:} services are delivered by registered and unregistered providers of the participant's choice
\end{itemize}

\section*{Clinical Features}
\begin{itemize}
    \item The NDIS funds supports, not health treatments, which remain the responsibility of the health system
    \item Participants manage their budgets with or without a support coordinator
    \item Plans are reviewed regularly, usually annually
    \item Supports must relate to the disability and represent value for money
    \item Participants can choose their providers and change them
    \item The NDIS works alongside Medicare, hospitals, and state disability services
\end{itemize}

\section*{Differential Diagnosis}
\begin{itemize}
    \item NDIS supports versus health system services, which are funded separately
    \item NDIS funding versus other government payments and concessions
    \item The roles of support coordinators, plan managers, and providers
\end{itemize}

\section*{When to Seek Medical Care}
People considering NDIS access should discuss their needs with their GP, who can provide evidence of disability and support the application. Seek medical care for health conditions as usual; the NDIS does not replace health care.

\textbf{Call 000 (or local emergency services) for:}
\begin{itemize}
    \item Medical emergencies affecting a person with disability
    \item Any situation where a person is at risk of harm
    \item Sudden deterioration in health
\end{itemize}

\section*{Diagnosis and Investigations}
\begin{itemize}
    \item \textbf{NDIS access request:} application with evidence of disability
    \item \textbf{Assessment:} functional assessment of support needs
    \item \textbf{Planning meeting:} discussion of goals and supports with a planner
    \item \textbf{Plan approval:} funding is approved for reasonable and necessary supports
    \item \textbf{Review:} plans are reviewed regularly and can be changed
\end{itemize}

\section*{Management}
\begin{itemize}
    \item \textbf{Access application:} applying through the NDIS with supporting evidence
    \item \textbf{Plan development:} setting goals and identifying supports
    \item \textbf{Choice of providers:} selecting therapy, support, and equipment providers
    \item \textbf{Support coordination:} help with implementing and managing the plan
    \item \textbf{Plan management:} financial management of the plan, by the participant or a plan manager
    \item \textbf{Using supports:} scheduling and using funded services
    \item \textbf{Review and appeals:} plan reviews and appeal processes for disagreements
\end{itemize}

\section*{Complications}
\begin{itemize}
    \item Complex application processes and delays
    \item Difficulties finding providers, especially in rural and remote areas
    \item Confusion about what the NDIS does and does not fund
    \item Gaps between NDIS supports and health services
    \item Financial and administrative stress for participants and families
    \item Disagreements about plan decisions
\end{itemize}

\section*{Prognosis and Outcomes}
For many people with disability, the NDIS has transformed access to supports, enabling greater independence, community participation, and quality of life. Outcomes depend on the quality of the plan and the availability of suitable providers. With good planning, coordination, and advocacy, most participants use their funding effectively and achieve their goals.

\section*{Prevention and Self-Care}
\begin{itemize}
    \item Learn about the NDIS before applying, including what it funds
    \item Gather medical evidence of your disability for the application
    \item Set clear goals for your plan
    \item Choose providers who suit your needs
    \item Use support coordination if the plan feels complex
    \item Attend plan reviews and request changes when needs change
    \item Seek advocacy support if you disagree with a decision
\end{itemize}

\section*{Sources and Further Reading}
The following reputable sources provide further information for healthcare students and patients seeking reliable, current advice about the NDIS.

\begin{itemize}
    \item National Disability Insurance Scheme. \textit{About the NDIS and how to access it.} https://www.ndis.gov.au/
    \item Disability Advocacy Network Australia. \textit{Advocacy and NDIS rights.} https://www.dana.org.au/
    \item Services Australia. \textit{Disability support information.} https://www.servicesaustralia.gov.au/
    \item Carer Gateway. \textit{Support for carers and families.} https://www.carergateway.gov.au/
    \item Healthdirect Australia. \textit{Disability services and the NDIS.} https://www.healthdirect.gov.au/
    \item People with Disability Australia. \textit{Information and advocacy for people with disability.} https://pwd.org.au/
\end{itemize}
